%------------------------------------------------------------------------------%
\documentclass[a4paper,12pt,fleqn]{article}

\usepackage{integracao_e_series}
\usepackage{systeme}

\ShowAnswers

%------------------------------------------------------------------------------%
\begin{document}

\makeHeader{Integrais e Séries}{Prova 2 -- Turma 8}{29/11/24}

% 01
%------------------------------------------------------------------------------%
\exercise{20}
Calcule a área entre as curvas
\(f(x) = -x e^x\),
\(g(x) = e^x\) e
\(x=1\)

\begin{answer}
  \begin{center}
    \includegraphics[width=0.8\linewidth]{area_entre_x_exp}
  \end{center}
  Encontrando o ponto onde os gráficos das funções $f$ e $g$ se interceptam
  \begin{align*}
    f(x)     & = g(x) \\
    -x e^{x} & = e^x  \\
    x        & = -1
  \end{align*}
  Como no intervalo \([-1, 1]\) a função $g$ é maior do que $f$ a área é dada
  pela integral
  \begin{align*}
    A & = \int_{-1}^1 g(x) - f(x) \,dx \\[2mm]
      & = \int_{-1}^1 e^x + x e^x \,dx \\[2mm]
      & = \int_{-1}^1 e^x \,dx  +  \int_{-1}^1x e^x \,dx
  \end{align*}
  Calculando a primitiva de \(h(x)=xe^x\) por integral por partes
  \(\int u\,dv = uv - \int v\,du\)
  \[
    H(x) = \int x e^x \,dx
  \]
  \[
    u = x
    \qquad
    du = dx
    \qquad
    dv = e^x\,dx
    \qquad
    v = e^x
  \]
  \begin{align*}
    H(x) & = xe^x - \int e^x\,dx \\[2mm]
         & = xe^x - e^x + C \\[2mm]
         & = (x - 1)e^x + C
  \end{align*}
  Portanto
  \begin{align*}
    A & = \int_{-1}^1 e^x \,dx  +  \int_{-1}^1x e^x \,dx \\[2mm]
      & = \left(e^x\right)\evalat{-1}{1} + \big((x - 1)e^x\big)\evalat{-1}{1} \\[2mm]
      & = e^1 -e^{-1} + \big((1 - 1)e^1\big) - \big((-1 - 1)e^{-1}\big) \\[2mm]
      & = e -\frac{1}{e} + 0 + \frac{2}{e} \\[2mm]
      & = e + \frac{1}{e}
  \end{align*}

\end{answer}

% 02
%------------------------------------------------------------------------------%
\exercise{20}
Use substituição simples para calcular a integral
\(
  \int \frac{a+bx^2}{\sqrt{3ax+bx^3}} dx
\)

\begin{answer}
  Escolhendo a substituição
  \[
    u = 3ax+bx^2
    \qquad
    du = \left(3a + 3bx^2\right) dx = 3\left(a + bx^2\right) dx
  \]
  portanto
  \[
    \left(a+bx^2\right)dx = \frac{du}{3}
  \]
  então
  \begin{align*}
    \int \frac{a+bx^2}{\sqrt{3ax+bx^3}} dx
    & = \int \frac{1}{\sqrt{3ax+bx^3}} \left(a+bx^2\right)dx \\
    & = \int \frac{1}{\sqrt{u}} \frac{du}{3} \\
    & = \frac{1}{3} \int u^{-\sfrac{1}{2}} du \\
    & = \frac{1}{3} \frac{u^{\sfrac{1}{2}}}{\sfrac{1}{2}} + c \\
    & = \frac{2}{3} \sqrt{u} + c \\
    & = \frac{2}{3} \sqrt{3ax+bx^2} + c
  \end{align*}
\end{answer}

% 03
%------------------------------------------------------------------------------%
\exercise{20}
Encontre a primitiva da função  \( f(x) = \cos^2(x)\sin(2x)  \)

\begin{answer}
  \begin{align*}
    F
    & = \int f(x) dx \\
    & = \int \cos^2(x)\sin(2x) dx \\
    & = \int \cos^2(x)2\sin(x)\cos(x) dx \\
    & = 2\int \cos^3(x)\sin(x) dx
  \end{align*}
  Fazendo a substituição
  \(
    \qquad
    u = \cos(x)
    \qquad
    du = -\sin(x)dx
  \)
  \begin{align*}
    F
    & = -2\int u^3 du \\
    & = -2 \frac{u^4}{4} + c \\
    & = -\frac{1}{2}\cos^4(x) + c
  \end{align*}
\end{answer}


% 04
%------------------------------------------------------------------------------%
\exercise{20}
Encontre a primitiva da função \(h(x) = \frac{x^5}{\sqrt{4-x^2}} \)

\begin{answer}
  \[
  H = \int \frac{x^5}{\sqrt{4-x^2}} dx
  \]
  Fazendo a substituição
  \(
  x = 2\sin(\theta)
  \qquad
  dx = 2\cos(\theta) d\theta
  \)
  \begin{spreadlines}{2ex}
    \begin{align*}
      H
      & = \int \frac{2^5 \sin^5(\theta)}{\sqrt{4-4\sin^2(\theta)}} 2\cos(\theta) d\theta \\
      & = 2^5 \int \frac{\sin^5(\theta)}{\sqrt{1-\sin^2(\theta)}} \cos(\theta) d\theta \\
      & = 2^5 \int \frac{\sin^5(\theta)}{\cos(\theta)} \cos(\theta) d\theta \\
      & = 2^5 \int \sin^5(\theta) d\theta \\
      & = 2^5 \int \left(\sin^2(\theta)\right)^2 \sin(\theta) d\theta \\
      & = 2^5 \int \left(1-\cos^2(\theta)\right)^2 \sin(\theta) d\theta
    \end{align*}
  \end{spreadlines}
  Fazendo a substituição
  \(
  \qquad
  u = \cos(\theta)
  \qquad
  du = -\sin(\theta)d\theta
  \)
  \begin{align*}
    H
    & = -2^5 \int \left(1-u^2\right)^2 du \\
    & = -2^5 \int 1-2u^2 + u^4\, du \\
    & = 2^5 \int 2u^2 - u^4 -1\, du \\
    & = 2^5 \left(\frac{2}{3}u^3 - \frac{1}{5}u^5 - u \right) + c \\
    & = 2^5 \left(\frac{2}{3}\cos^3(\theta)
    - \frac{1}{5}\cos^5(\theta)
    - \cos(\theta) \right) + c
  \end{align*}
  Para calcular \(\cos(\theta)\) usamos que \(\sin(\theta) = \sfrac{x}{2}\), portanto
  a hipotenusa é 2 e o cateto oposto é $x$ assim o cateto adjacente é
  \(
  a = \sqrt{2^2-x^2} = \sqrt{4-x^2}
  \)
  e o cosseno é
  \[
  \cos(\theta) = \frac{\sqrt{4-x^2}}{2}
  \]
  Voltando para a integral
  \begin{spreadlines}{2ex}
    \begin{align*}
      H
      & = 2^5\left[
      \frac{2}{3}\left(\frac{\sqrt{4-x^2}}{2}\right)^3
      - \frac{1}{5}\left(\frac{\sqrt{4-x^2}}{2}\right)^5
      - \frac{\sqrt{4-x^2}}{2}
      \right] + c \\
      & = 2^5\left[
      \frac{2}{3\cdot 2^3}\left(\sqrt{4-x^2}\right)^3
      - \frac{1}{5\cdot 2^5}\left(\sqrt{4-x^2}\right)^5
      - \frac{1}{2} \sqrt{4-x^2}
      \right] + c \\
      & = \sqrt{4-x^2}
      \left[
      \frac{2^3}{3}\left(\sqrt{4-x^2}\right)^2
      - \frac{1}  {5}\left(\sqrt{4-x^2}\right)^4
      - 2^4
      \right] + c \\
      & = \sqrt{4-x^2}
      \left[
      \frac{8}{3}\left(4-x^2\right)
      - \frac{1}{5}\left(4-x^2\right)^2
      - 16
      \right] + c \\
      & = \sqrt{4-x^2}
      \left[
      \frac{32-8x^2}{3}
      - \frac{16-8x^2+x^4}{5}
      - 16
      \right] + c \\
      & = \sqrt{4-x^2}
      \left[
      \frac{5\cdot 32 - 5\cdot 8x^2}{15}
      -\frac{3\cdot 16 - 3\cdot 8x^2 + 3x^4}{15} - \frac{16\cdot 15}{15}
      \right] + c \\
      & = \frac{\sqrt{4-x^2}}{15}
      \Big[
      16\cdot 5\cdot 2 - 16 \cdot 3 - 16\cdot 5 \cdot 3
      - 5\cdot 8x^2 + 3\cdot 8x^2
      - 3x^4
      \Big] + c \\
      & = \frac{\sqrt{4-x^2}}{15}
      \Big[
      - 16\cdot 5 -16 \cdot 3
      - 2\cdot 8x^2
      - 3x^4
      \Big] + c \\
      & = -\frac{\sqrt{4-x^2}}{15}
      \Big[
      16\cdot 8
      + 2\cdot 8x^2
      + 3x^4
      \Big] + c \\
      & = -\frac{\sqrt{4-x^2}}{15}
      \Big[ 128 + 16x^2 + 3x^4 \Big] + c
    \end{align*}
  \end{spreadlines}
\end{answer}

% 05
%------------------------------------------------------------------------------%
\exercise{20}
Calcule o volume do sólido de revolução gerado pela rotação,
em torno da reta \(x=-1\), da região contida entre as curvas
\[
  f(x) = \frac{1}{4x-x^2} \qquad
  y = 0 \qquad
  x = 1 \qquad
  x = 3
\]

\begin{answer}
  Volume por castas cilíndricas
  \begin{align*}
    V & = \int_a^b 2\pi r(x) h(x) dx \\
    a & = 1 \\
    b & = 3 \\
    r & = x+1 \\
    h & = \frac{1}{4x-x^2}
  \end{align*}
  portanto
  \[
    V = 2\pi \int_1^3 \frac{x+1}{4x-x^2} dx
  \]
  Calculando a integral indefinida
  \[
    F = \int \frac{x+1}{4x-x^2} dx = \int \frac{x+1}{x(4-x)} dx
  \]
  Por frações parciais temos
  \begin{align*}
    \frac{x+1}{x(4-x)} & = \frac{A}{x} + \frac{B}{4-x} \\
    x+1                & = A(4-x) + Bx                 \\
                       & = (B-A)x + 4A
  \end{align*}
  Igualando os coeficientes
  \[
    \systeme[AB]{B-A=1,4A=1}
  \]
  obtemos os valores
  \(A=\frac{1}{4}\) e
  \(B=\frac{5}{4}\)
  portanto
  \begin{align*}
    F & = \int \frac{1}{4x} + \frac{5}{4(4-x)} dx \\
      & = \frac{1}{4}\ln(x) - \frac{5}{4}\ln(4-x) + c
  \end{align*}
  Voltando ao volume temos
  \begin{align*}
    V & = 2\pi F(x)\evalat{1}{3} \\
      & = 2\pi \left[\frac{1}{4}\ln(x) - \frac{5}{4}\ln(4-x)\right] \evalat{1}{3} \\
      & = 2\pi \left[
                 \left(\frac{1}{4}\ln(3) - \frac{5}{4}\ln(4-3)\right)
               - \left(\frac{1}{4}\ln(1) - \frac{5}{4}\ln(4-1)\right)
               \right] \\
      & = 2\pi \left(\frac{1}{4}\ln(3) + \frac{5}{4}\ln(3)\right) \\
      & = 2\pi \frac{6}{4}\ln(3) \\
      & = 3\pi \ln(3)
  \end{align*}
\end{answer}

%------------------------------------------------------------------------------%
\end{document}
%------------------------------------------------------------------------------%
